\documentclass[preview]{standalone}
\usepackage[english]{babel}
\usepackage{amsmath}
\usepackage{amssymb}
\usepackage{ragged2e}
\setlength\parindent{0pt}
\begin{document}
\begin{center}
\special{dvisvgm:raw <g id='unique000'>}\parbox{9cm}{\noindent{}\justifying{\special{dvisvgm:raw </g>}\special{dvisvgm:raw <g id='unique001'>}\begin{itemize} \item A group is uniscalar if the scale of every element is 1.\end{itemize}\special{dvisvgm:raw </g>}\special{dvisvgm:raw <g id='unique002'>}\begin{itemize} \item A uniscalar group is ``very close'' to being discrete without necessarily being discrete.\end{itemize}\special{dvisvgm:raw </g>}\special{dvisvgm:raw <g id='unique003'>}\begin{itemize} \item For (P)-closed groups with compact-open vertex stabilisers, we use our formula to show when they are uniscalar.\end{itemize}\special{dvisvgm:raw </g>}\special{dvisvgm:raw <g id='unique004'>}}}\special{dvisvgm:raw </g>}
\end{center}
\end{document}